\documentclass[10pt,a4paper]{article}
\usepackage[margin=1.6cm,top=1.3cm,bottom=1.3cm]{geometry}
\usepackage[T1]{fontenc}
\usepackage{lmodern}
\usepackage{amsmath}
\usepackage{enumitem}
\usepackage{titlesec}
\usepackage[hidelinks]{hyperref}
\usepackage{xcolor}
\newcommand{\todo}[1]{\textcolor{red}{[#1]}}

\titleformat{\section}{\normalsize\bfseries}{\thesection.}{0.4em}{}
\titlespacing*{\section}{0pt}{5pt}{2pt}
\setlist{nosep,leftmargin=1.2em}
\setlength{\parindent}{0pt}
\setlength{\parskip}{2pt}
\pagestyle{empty}

\begin{document}

\begin{center}
{\large\bfseries Can Next Week's Dengue Surge Be Flagged Early? An Early-Warning Model for Singapore}\\[2pt]
{\small Project 1.1 (public health, intermediate) \quad$\cdot$\quad Statistical Learning group project
\quad$\cdot$\quad \todo{Group members}}
\end{center}

\section{Question and motivation}
Dengue is endemic in Singapore and arrives in irregular epidemic waves (e.g.\ 2013--14, 2019--20, 2022).
\textbf{Primary question:} can recent dengue counts and lagged weather (rainfall, temperature, humidity)
identify weeks with a materially elevated risk of a surge in the following week(s)? An early, calibrated
risk signal can support inspection, communication and vector-control planning, even though the model
does not establish what causes a surge. We frame the task as \emph{prediction}, not causal inference.

\section{Data and curation}
\begin{itemize}
  \item \textbf{Cases:} MOH weekly dengue fever (DF) and dengue haemorrhagic fever (DHF), data.gov.sg, 2014--2018
        (main analysis); MOH Weekly Infectious Disease Bulletin 2012--2022 for an extended robustness analysis
        covering more outbreaks. DF is the target series; DHF is reported separately ($<$1 case/week on average).
  \item \textbf{Weather:} NEA/MSS daily records (Changi), aggregated to epidemiological weeks and lagged 0--12 weeks;
        reanalysis (Open-Meteo) only to fill gaps, cross-checked against NEA monthly series.
  \item \textbf{Curation:} a complete epi-week calendar (Sun--Sat, MMWR numbering, incl.\ 53-week years);
        datasets joined on overlapping weeks only; every predictor uses information available at forecast week $t$.
\end{itemize}

\section{Surge definitions (fixed before any model is fitted)}
For horizon $h\in\{1,2,4\}$ weeks ($h=1$ is primary) and DF counts $y_t$:
\begin{itemize}
  \item \textbf{Seasonal percentile (S):} $y_{t+h} > Q_{0.90}$ of counts in the same season ($\pm$2.5 weeks) of earlier years.
  \item \textbf{Historical threshold (A):} $y_{t+h} > \bar y_{t-51:t} + 2\,s_{t-51:t}$ (past 52 weeks).
  \item \textbf{Week-on-week rise (B):} $y_{t+h} \ge (1+p_h)\,y_t$ and $y_{t+h}\ge 50$.
  \item \textbf{Onset variants of A and S:} only weeks where $y_t$ is \emph{not yet} above the threshold, so that
        the model is judged on genuine early warnings rather than on the persistence of an ongoing epidemic.
\end{itemize}
Thresholds are tuned once so that 10--20\% of weeks are surges (preliminary: A 19\%, S 19\%, B 12--18\%).

\section{Methods}
\begin{itemize}
  \item \textbf{Baselines:} base rate; \emph{seasonal} (historical surge rate for that week of the year);
        \emph{persistence} (current distance to the threshold, or last week's growth).
  \item \textbf{Classification:} logistic regression with lagged cases and/or weather; $\ell_1$-regularised
        logistic regression (penalty by forward-chaining CV) for variable selection; CART (depth 3, interpretable
        rules); random forest.
  \item \textbf{Count regression:} Poisson GLM and ridge regression on $\log(1+y)$, converted to $P(\text{surge})$
        through a predictive distribution so that all models are compared on the same scale.
\end{itemize}

\section{Validation and evaluation}
Rolling-origin evaluation \emph{by week} with an expanding window: at each forecast week the model is trained only
on weeks whose outcome is already known (purged by $h$), with no future weather or later case revisions in the
predictors. We report \textbf{PR-AUC} (main metric, surges are rare), recall, precision and alarms per year at an
alarm threshold set by the cost ratio of a missed surge to a false alarm, \textbf{calibration} (Brier score,
reliability diagrams), 95\% block-bootstrap intervals, and the gain over the persistence and seasonal baselines.

\section{Anticipated limitations (defence focus)}
\begin{itemize}
  \item \textbf{Reporting and revisions:} only final counts are public; sensitivity run assuming a 1-week reporting delay.
  \item \textbf{Immunity and serotype shifts} (e.g.\ DENV-3 in 2020) change the case--weather relation; performance
        is reported per test year.
  \item \textbf{Few large outbreaks:} small effective sample; simple, regularised models and bootstrap intervals.
  \item \textbf{Cost of false alarms:} alarm rules compared across cost ratios (1:1 to 9:1), expressed as alarms per year.
\end{itemize}

\section{Deliverables and timeline}
Wk 1--2: data curation and surge definitions \;$\cdot$\; Wk 3--4: baselines, logistic and count models
\;$\cdot$\; Wk 5: trees, regularisation, model selection \;$\cdot$\; Wk 6: robustness and defence analyses
\;$\cdot$\; Wk 7--8: report and reproducible code (public GitHub pipeline).

\vspace{2pt}
{\footnotesize \textbf{References.}
Hii~Y.L.\ et al.\ (2012), Forecast of dengue incidence using temperature and rainfall, \emph{PLoS NTD} 6(11):e1908.
Shi~Y.\ et al.\ (2016), Three-month real-time dengue forecast models: an early warning system for outbreak alerts
and policy decision support in Singapore, \emph{Environ.\ Health Perspect.} 124(9):1369--1375.\par}

\end{document}
